% =====================================================================
%  04-recursos.tex - Recursos procesales
% =====================================================================

\section{Recursos procesales}\label{sec:impl-recursos}

Esta sección implementa la revocación y la apelación que el modelo jurídico (\cref{sec:modelo-juridico}) deja para este capítulo; no tiene un RF propio y aporta a RF-07 y RF-08 (\cref{tab:rf-calendario}) mediante sus audiencias propias y la creación contextual de plazos (marco normativo en la \cref{tab:terminos-fatales}). La \cref{fig:impl-rec-panorama} muestra sus cuatro piezas; los avisos se tratan en la \cref{sec:impl-alertas-audiencias-recursos}. En todas, Owner y el Litigante asignado escriben, el Paralegal asignado consulta y Cliente queda excluido; cada lectura confirma su auditoría antes de entregar datos, y el cierre administrativo conserva consultas y repetición exacta autorizada, pero bloquea mutaciones.

\begin{figure}[!htbp]
  \centering
  \begin{tikzpicture}[
  fuente/.style={caja, text width=3.35cm, minimum height=0.9cm},
  act/.style={caja, text width=3.75cm, minimum height=0.95cm}]

  % Zonas
  \draw[zona] (0,0) rectangle (4.05,-4.45);
  \draw[zona] (11.45,0) rectangle (16,-4.45);
  \node[zonaLabel] at (0,0) {FUENTES HISTÓRICAS};
  \node[zonaLabel] at (11.45,0) {ACTIVIDADES DEL EXPEDIENTE};

  % Fuentes
  \node[fuente] (res) at (2.025,-1.0) {\textbf{Resolución impugnada}\\ revisión exacta};
  \node[fuente] (rec) at (2.025,-2.25) {\textbf{Recurrentes}\\ hasta 32, ficha opcional};
  \node[fuente] (sop) at (2.025,-3.5) {\textbf{Soportes documentales}\\ versión y huella exactas};

  % Recurso y actos
  \node[cajaOsc, text width=4.4cm, minimum height=1.25cm] (rcs) at (7.75,-1.2)
    {\textbf{Recurso declarado}\\ revocación o apelación\\ raíz y revisiones inmutables};
  \node[cajaClara, text width=4.4cm, minimum height=1.15cm] (actos) at (7.75,-3.5)
    {\textbf{Actos declarados}: interposición, admisión, inadmisibilidad, desistimiento y resolución};
  \draw[flecha] (actos) -- node[etq, right=2pt, align=left]{cada acto añade\\ una revisión} (rcs);
  \draw[flecha] (res.east) -- (res.east -| rcs.west);
  \draw[flecha] (rec.east) -- (4.35,-2.25) |- (rcs.west |- 0,-1.4);
  \draw[flecha] (sop.east) -- (actos.west);
  \draw[flecha] (sop.east) -- (4.8,-3.5) |- (rcs.west |- 0,-1.7);

  % Actividades
  \node[act] (exi) at (13.725,-0.85) {\textbf{Audiencia ordinaria o plazo} que ya existen};
  \node[act] (aud) at (13.725,-2.2) {\textbf{Audiencia propia}: programación con su asociación inicial};
  \node[act] (pla) at (13.725,-3.55) {\textbf{Plazo} creado con su asociación en una confirmación};
  \draw[flecha] (rcs.east |- 0,-0.66) -- node[etq, above]{asociación} (exi.west |- 0,-0.66);
  \draw[flechaSuave] (exi.west |- 0,-1.08) -- node[etq, below, align=center]{consulta\\ inversa} (rcs.east |- 0,-1.08);
  \draw[flecha] (rcs.east |- 0,-1.72) -- (10.85,-1.72) |- (aud.west);
  \draw[flecha] (10.85,-2.2) |- (pla.west);
\end{tikzpicture}

  \caption{Panorama de los recursos procesales: fuentes históricas exactas y tres caminos hacia actividades, ninguno con efectos jurídicos automáticos.}
  \label{fig:impl-rec-panorama}
\end{figure}

\subsection{Recursos procesales declarados}\label{sec:impl-recursos-declarados}

El registro de recursos representa la revocación y la apelación como historias independientes de la etapa ordinaria del expediente. Su finalidad es conservar lo declarado por el usuario, la resolución seleccionada y los soportes que acompañan cada captura. No introduce una cuarta etapa ni interpreta que registrar un recurso acredite su presentación, admisión, firmeza o efectos sobre otros actos. El contrato implementado está en \rutarepo{docs/procedural-resources-api.md}; la separación entre registro organizativo y efectos procesales se fundamenta en \rutarepo{docs/adr/0040-resources-linked-to-historical-resolutions.md}.

\subsubsection{Identidad, declaración y fuentes históricas}

Cada recurso tiene una identidad y una secuencia propia de revisiones. El registro distingue tipo de recurso, modalidad oral o escrita, título organizativo, parte impugnada, motivos, autoridades y tiempos declarados. Una declaración conocida conserva su valor; una declaración desconocida conserva el motivo de ese desconocimiento. En los campos opcionales, la ausencia se mantiene distinta de una declaración desconocida. Una fecha sin hora no se completa con medianoche ni se convierte en un instante preciso.

La resolución se selecciona mediante identidad y revisión exactas dentro del mismo expediente. La preparación verifica sus recibos y conserva una proyección legible de esa revisión. Por ejemplo, elegir la revisión 1 no implica reemplazarla por la revisión 2 al confirmar, aunque ésta sea la cabeza consultable. Pueden coexistir varios recursos referidos a la misma resolución. La selección histórica tampoco atribuye vigencia o elegibilidad jurídica a la fuente.

Las personas recurrentes conservan nombre y papel declarados. Su vínculo opcional con el directorio identifica una revisión histórica verificada y no las convierte en cuentas de acceso. El límite de 32 capturas por recurso es técnico. Se rechaza la repetición de una misma ficha del directorio; dos personas declaradas sin ficha pueden compartir nombre. La resolución impugnada exige un soporte documental exacto, identificado por documento, versión, digest y localizador.

\subsubsection{Actos e historia organizativa}

Los actos declarados distinguen interposición, admisión, inadmisibilidad, desistimiento y resolución. Cada acto tiene identidad, revisiones, modalidad, tiempo, autoridad, declaración y evidencia propios. Registrar o corregir un acto añade una revisión al recurso; la corrección referencia además la captura anterior de ese mismo acto. Por ello, la revisión anterior del acto puede encontrarse varias revisiones atrás en la historia del recurso.

Cada acto admite entre uno y dos localizadores de evidencia, conforme al límite técnico del lote documental existente. Dos localizadores de una misma versión conservan sus funciones y comparten la admisión; un digest contradictorio para esa versión provoca rechazo. Se reutiliza la admisión PDF/DOCX, sin afirmar que esos formatos determinen la suficiencia jurídica de la evidencia. Una constancia documental puede respaldar un acto declarado oral sin convertirlo en un acto escrito.

El archivo y la reactivación son cambios organizativos que preservan valores y fuentes. Archivar no registra un desistimiento. La cabeza del recurso expone el acto que se haya registrado en esa revisión, no una lista implícita de todos los actos anteriores. La historia permite consultar cada captura y reconstruir las revisiones de una identidad de acto.

\subsubsection{Preparación, confirmación y persistencia}

El servicio separa seis intenciones: registrar y corregir recurso, registrar y corregir acto, archivar y reactivar. La preparación reúne fuentes exactas, verifica las nuevas versiones documentales en un lote acotado y genera un digest del envío que el usuario revisa. No reserva la operación ni mantiene bloqueadas las fuentes mientras se muestra el formulario. Los soportes admitidos que se retienen se contrastan con su captura, sin sustituirlos por versiones más recientes ni readmitirlos silenciosamente.

La confirmación vuelve a autenticar al principal completo y contrasta la revisión esperada del recurso, la revisión del acto cuando corresponde, la evidencia preparada y las observaciones administrativas y de etapa. El adaptador PostgreSQL revalida cuenta activa, rol y asignación bajo el bloqueo de auditoría compartido. Añade la revisión, el recibo y el evento de auditoría dentro de la misma transacción. Si el contexto observado cambió, devuelve un conflicto; no confirma sobre una nueva base sin una decisión explícita.

El recibo vincula comando, autor con correo capturado, valores, fuentes, contexto observado y digest de la captura anterior. La captura final añade el instante UTC de confirmación. Las codificaciones canónicas \texttt{PRSC1}, \texttt{PRAC1}, \texttt{PRSS1}, \texttt{PRTX1} y \texttt{PRCP1} (\cref{tab:anexo-canones}) distinguen, respectivamente, valores del recurso, valores del acto, fuentes, envío y captura; su comprobación demuestra coherencia técnica y no aplicabilidad jurídica. Una repetición con la misma identidad de operación, actor y comando conserva el recibo original. Un uso diferente de esa identidad produce conflicto. La autorización vigente también se exige para recuperar una operación ya confirmada.

La migración \rutarepo{migrations/0022_procedural_resources.sql} (\cref{tab:anexo-migraciones}) incorpora raíces, actos y revisiones inmutables. Sus guardas y el catálogo de arranque controlan secuencias, relaciones y privilegios de inserción. El inventario utiliza el lector histórico del adaptador, recorre revisiones en lotes acotados y rechaza capturas incoherentes sin escribir nuevos eventos. La lectura de una corrección también comprueba que la raíz del acto conserve su captura inicial exacta. La \cref{sec:pruebas-recursos-declarados} reúne las pruebas del adaptador, incluida su exportación y restauración, y la aceptación integrada de API y navegador.

\subsubsection{Consulta autorizada e interacción en Qadra}

Owner puede gestionar los recursos; Litigante requiere asignación vigente al expediente. Paralegal asignado consulta y Cliente conserva la denegación. El cierre administrativo permite consultar la historia y bloquea nuevas mutaciones. La lectura devuelve datos sólo después de confirmar su auditoría y de la reautenticación del servicio. Las capturas conservan el correo del autor original aunque su cuenta cambie posteriormente.

La API ofrece colección, cabeza, revisión exacta e historia bajo la ruta del expediente. El listado selecciona cabezas antes de aplicar los filtros de tipo y estado, y utiliza continuación exclusiva por identidad. Admite hasta 100 recursos por página, con 20 por omisión; la historia descendente admite hasta 20 revisiones, con 10 por omisión. Estos límites de consulta no son límites procesales. El transporte rechaza campos desconocidos, cuerpos ambiguos y discordancias entre ruta, recurso, acto e intención.

Qadra presenta formulario, revisión previa a confirmar, detalle e historia. Conflictos, respuestas inciertas y revocación siguen los mecanismos comunes de las \cref{sec:impl-qadra-borradores,sec:impl-qadra-conciliacion}: el borrador se conserva hasta aceptar expresamente la base actual, la conciliación consulta el recibo exacto de la operación sin tomar una ausencia temporal como prueba de fracaso, y una consulta nuevamente denegada retira los datos visibles.

\subsubsection{Frontera pendiente del flujo completo}

El registro y la consulta de recursos declarados se complementan con las demás piezas de esta sección: los vínculos con audiencias y plazos existentes (\cref{sec:impl-asociaciones-recursos}), la lectura inversa que recorre recursos desde una actividad o alerta y regresa a su origen (\cref{sec:impl-recursos-desde-actividad}), las audiencias propias de un recurso (\cref{sec:impl-audiencias-recursos}) y la creación conjunta de un plazo con su vínculo (\cref{sec:impl-plazo-contextual}). Ese recorrido de lectura no acredita causalidad jurídica del aviso. Ninguno de los comandos de registro avanza la etapa, suspende términos o determina elegibilidad. El sistema no activa términos automáticamente a partir de un acto declarado: cada plazo se registra mediante una instrucción explícita, y la cobertura jurídica aceptada se limita al corpus finito de ocho supuestos verificado en la \cref{sec:pruebas-corpus-juridico}. Los criterios de aceptación del alcance completo de recursos se conservan en \rutarepo{docs/procedural-resources-scope.md}; el registro declarado no los sustituye por fechas manuales ni acredita por sí solo que el flujo completo de recursos esté concluido.


\subsection{Asociaciones de recursos con actividades existentes}\label{sec:impl-asociaciones-recursos}

La asociación organizativa permite relacionar un recurso declarado, y opcionalmente uno de sus actos, con una audiencia o un plazo que ya existen en el mismo expediente. El vínculo tiene identidad e historia independientes de sus extremos. Su registro no programa otra audiencia, recalcula un término ni genera una segunda alerta. La decisión se documenta en \rutarepo{docs/adr/0041-exact-resource-activity-associations.md} y el intercambio HTTP en \rutarepo{docs/resource-activities-api.md}.

\subsubsection{Selección exacta y cabeza esperada}

La selección conserva identidad, revisión y digest del recurso. Cuando incluye un acto, identifica también la revisión del recurso que contiene esa captura. Ambas selecciones son independientes: puede vincularse la revisión 1 del recurso junto con un acto conservado en la revisión 2, aunque una corrección posterior sea la cabeza actual. El servidor verifica que cada fuente pertenezca al expediente y al recurso solicitados; no sustituye la selección por la última revisión ni vuelve a admitir el contenido documental retenido.

La audiencia ordinaria se identifica por su revisión y digest de envío; el plazo y la audiencia propia de un recurso, por su revisión y digest de captura. La cabeza esperada del recurso se declara aparte de esas referencias históricas. Preparar reúne las fuentes verificadas y el contexto administrativo; confirmar vuelve a contrastarlos bajo el bloqueo compartido de auditoría. Una revisión concurrente produce conflicto y exige aceptar expresamente otra base.

La vinculación crea la revisión 1. La desvinculación añade una revisión 2 con motivo, preservando selección y fuentes. Es terminal para esa identidad; un vínculo posterior tiene otra raíz. Pueden existir varias asociaciones con una misma actividad sin duplicar la actividad ni sus avisos. El archivo del recurso impide crear vínculos nuevos, pero permite desvincular los existentes cuando el expediente está activo.

\subsubsection{Historia y estado actual separados}

Detalle y lista exponen dos partes: la asociación histórica con sus capturas, y la cabeza actual verificada de la actividad. Ambas se resuelven en una transacción autorizada. Una audiencia histórica puede conservar su programación original aunque la cabeza actual esté cancelada. Para un plazo, la evaluación capturada permanece legible, pero sólo la proyección operativa de la cabeza actual determina si se muestra un vencimiento vigente. Un plazo retirado, pendiente o afectado por cambios no recupera una fecha operativa desde su cálculo histórico.

La consulta toma un instante UTC común después de adquirir el bloqueo. Todas las filas de la página y las proyecciones operativas de plazo comparten ese corte. El servicio comprueba que esté entre el inicio de la solicitud y el retorno del almacén. Así admite una cabeza confirmada mientras el lector esperaba, sin aceptar una observación anterior a la solicitud o posterior al retorno. La historia y los recibos de escritura no incorporan esa afirmación de vigencia.

\subsubsection{Persistencia, permisos e interfaz}

Las migraciones de asociaciones, conservadas en el directorio \rutarepo{migrations/}, incorporan dos tablas: raíces y revisiones de asociación. La transacción añade revisión, recibo y auditoría conjuntamente. El catálogo estricto comprueba estructura y privilegios; el inventario reconstruye capturas históricas en lotes acotados, sin modificar datos al abrir. El respaldo conserva ambas tablas y sus dependencias, no sólo los identificadores de los extremos. La \cref{sec:pruebas-asociaciones-recursos} reporta la aceptación API con restauración, que comparó capturas y respuestas y validó por separado los instantes de consulta nuevos, junto con los recorridos de navegador real.

Owner gestiona todos los expedientes; Litigante asignado gestiona y Paralegal asignado consulta. Cliente no accede. El servicio reautentica al principal completo y el adaptador comprueba cuenta, rol y asignación vigentes. El cierre conserva lecturas y repetición exacta autorizada de una operación confirmada, pero impide nuevas mutaciones.

Qadra presenta selectores de actividad, revisión y acto opcional, revisión previa a confirmar, listado filtrado, detalle e historia; la captura vinculada y la actividad actual ocupan áreas separadas. Conflictos, respuestas inciertas y revocación siguen los mecanismos comunes de las \cref{sec:impl-qadra-borradores,sec:impl-qadra-conciliacion}: un conflicto conserva selección y motivo hasta aceptar una nueva base, ninguna coincidencia visual sustituye al recibo exacto, y el reenvío explícito tras un resultado incierto conserva el mismo comando, operación y digest, sin preparar otra identidad.

Las asociaciones sólo vinculan actividades que ya existen. La creación conjunta de un plazo y su asociación tiene el contrato independiente de la \cref{sec:impl-plazo-contextual}, y la programación de una audiencia propia crea su asociación inicial en la misma transacción (\cref{sec:impl-audiencias-recursos}). La consulta inversa de la \cref{sec:impl-recursos-desde-actividad} completa un recorrido de lectura desde actividades y avisos, sin crear actividades adicionales. Los vínculos no activan términos ni amplían la cobertura jurídica delimitada al final de la \cref{sec:impl-recursos-declarados}, y no modifican los objetivos aprobados.

\subsubsection{Consulta inversa y regreso a la actividad}\label{sec:impl-recursos-desde-actividad}

La consulta inversa parte de una audiencia ordinaria o de un plazo identificado dentro del expediente y devuelve las asociaciones actuales de sus recursos. Utiliza las rutas de audiencia y plazo con sufijo \texttt{resource-associations}, sin recorrer en el navegador todos los recursos del expediente. La lista muestra vínculos activos de manera predeterminada, permite consultar los desvinculados o ambos estados y pagina por UUID ascendente con un límite de 1 a 100. El estado se filtra después de seleccionar la cabeza de cada asociación y antes de aplicar el límite. Cada continuación conserva expediente, destino y filtro; las páginas son observaciones independientes, no una instantánea histórica retenida entre peticiones.

Cada respuesta incluye un instante UTC común incluso cuando no contiene filas. La transacción comprueba el destino, las capturas de recurso, acto y actividad, sus cabezas actuales y la auditoría de lectura; el servicio reautentica al principal antes de entregar datos. Owner consulta todos los expedientes y Litigante o Paralegal necesitan asignación vigente. Cliente permanece denegado. El cierre del expediente y el archivo del recurso conservan las lecturas autorizadas. Un error de acceso no se representa como ausencia de relaciones.

Qadra presenta el panel \enquote{Recursos relacionados} en el detalle de actividad. Se distinguen la revisión abierta, la revisión de actividad capturada por cada vínculo y la cabeza actual observada. La lista representa el estado actual de las asociaciones: una relación registrada después de emitir una alerta puede aparecer, sin afirmar que existiera entonces ni que causara el aviso. La selección conserva la identidad, revisión y digest tanto del recurso como de la asociación, y su apertura comprueba las capturas exactas. Consultar el recurso actual sigue siendo una acción diferente.

El recorrido conserva el regreso a la revisión de actividad desde la que se abrió el recurso. Si comenzó en una alerta, mantiene también el regreso a la bandeja y sus filtros. Las consultas no marcan el aviso como leído ni modifican su origen; la revocación y las respuestas tardías siguen los mecanismos comunes de las \cref{sec:impl-qadra-borradores,sec:impl-qadra-conciliacion}. El contrato está en \rutarepo{docs/activity-resource-links-api.md} y la decisión en \rutarepo{docs/adr/0059-activity-resource-navigation.md}. Su verificación se distingue de la evidencia de las asociaciones en la \cref{sec:pruebas-recursos-desde-actividad}.


\subsection{Audiencias propias de recursos}\label{sec:impl-audiencias-recursos}

Un recurso escrito puede originar una audiencia propia: alegatos de apelación o audiencia de revocación escrita por complejidad (artículos 466 y 476--477 del CNPP \cite{cnpp_consolidado2025}). La clasificación no declara procedencia, no deriva una audiencia separada de una revocación oral ni exige etapa ordinaria; es una familia separada de las audiencias ordinarias y de su formato canónico (\cref{sec:impl-audiencias}). Captura tipo, instante con su desfase original, modalidad, lugar, nota opcional, hasta 32 participantes por ficha y revisión (sin repetir, ordenados por UUID) y una base declarada con soporte exacto (\texttt{RHEAR1}, \cref{tab:anexo-canones}). No completa la fecha ni calcula el señalamiento desde una norma: la precisión debe bastar para el instante comunicado.

La confirmación sigue la \cref{fig:impl-audiencias-recibo} con las huellas \texttt{RHPR1} (revisión preparada) y \texttt{RHCR1} (captura); el soporte debe ser uno ya admitido en el recurso o el acto, y una sola transacción escribe audiencia R1, asociación inicial y origen. La huella de captura liga cada participante a su identidad y revisión, y la captura rechaza tiempos anteriores a sus fuentes o posteriores al reloj. Dos envíos simultáneos de la misma operación recupera la creación original; una repetición explícita conserva comando y huella y, si la operación no existe, vuelve a exigir las condiciones de creación. La asociación inicial marca que la audiencia pertenece al recurso que la creó y no puede vincularse a otro. El inventario exige asociación y origen para cada una, de modo que una restauración incompleta no habilita una creación duplicada (\cref{tab:anexo-migraciones}).

Desvincular, archivar o cerrar no cancela el señalamiento, no impide consultarlo ni altera origen, captura o asociación inicial. En la agenda tiene rango propio del cursor y, a igual instante, sigue el orden de familias de la \cref{sec:implementacion-agenda-combinada}. Las consultas usan páginas de hasta veinte resultados con cursor exclusivo y sin instantánea común. Qadra abre el detalle histórico en Agenda tras verificar el expediente y coteja la revisión exacta sin deducir la actual (\cref{sec:impl-qadra-conciliacion,sec:pruebas-audiencias-recursos}).

\subsection{Creación contextual de plazos en recursos}\label{sec:impl-plazo-contextual}

Owner o el Litigante asignado crean desde el recurso un plazo y su vínculo con una sola confirmación, con perfil, fuente temporal, calendario, responsable y políticas del registro ordinario; el recurso aporta contexto, no fuente temporal ni aplicabilidad, y un cálculo bloqueado conserva sus motivos sin vencimiento manual. Una transacción escribe plazo, asociación y eventos, ligados por el recibo compuesto \texttt{RDLTX1} y un marcador auditado (\cref{tab:anexo-canones,anexo:plazo-contextual}). El plazo sigue el seguimiento ordinario y entra en la agenda si tiene vencimiento operativo; corregirlo no reescribe el vínculo y desvincularlo no lo retira; tampoco programa audiencias del recurso.

\paragraph{Límites.} Ningún comando de recursos avanza la etapa, suspende términos ni determina elegibilidad, y ningún acto declarado activa términos: cada plazo exige una instrucción explícita. La cobertura jurídica se limita al corpus de ocho supuestos (\cref{sec:pruebas-corpus-juridico}) y el alcance completo sigue pendiente según \rutarepo{docs/procedural-resources-scope.md}.

\Fuentes{documentos \texttt{procedural-resources-api}, \texttt{resource-activities-api}, \texttt{activity-resource-links-api}, \texttt{resource-hearings} y \texttt{resource-deadlines-api} de \rutarepo{docs/}; ADR 0040, 0041, 0057, 0059 y 0069 de \rutarepo{docs/adr/}.}
